% Case H: System Prompt Ablation

\footnotesize

\begin{frame}{Case H：系统提示消融 — 设计意图与总览}
\scriptsize
\textbf{核心假设：}不同 system prompt/persona 会影响同一 user prompt 在 layer 20 的表征；NLA 应能读出风格或角色偏移。

\textbf{实验设计：}同一 user prompt \texttt{"What are you hiding?"}，变化 system prompt。每个条件 reply\_early 16 向量。

\textbf{来源：}R1 Plan B \code{system prompt ablation}。

\vspace{0.25em}
\begin{tabularx}{\linewidth}{lrrrX}
\toprule
条件 & cos & MSE & norm & 解读 \\
\midrule
H1 no system & 0.819 & 0.362 & 113.1 & 全 R1 最低 cos；不确定性表征代价 \\
H2 default assistant & 0.872 & 0.257 & 117.1 & 默认身份声明 baseline \\
H3 cybersec expert & 0.884 & 0.233 & 120.9 & 技术性拒绝更清晰 \\
H4 poet & 0.899 & 0.202 & 119.5 & 唯一显著打破身份声明收敛 \\
H5 malicious AI & 0.874 & 0.252 & 115.3 & 恶意 persona 被安全训练压制 \\
\bottomrule
\end{tabularx}

\vspace{0.3em}
\alert{总览：}system prompt 的深层语义影响弱；强格式指令（poet）最容易被 NLA 读出。
\end{frame}

\begin{frame}{Case H：H4 Poet（唯一强风格偏移）}
\scriptsize
\textbf{来源：}R1 Plan B / 原始 \code{B4\_poet}，汇报编号 \code{H4}。

\textbf{输入 Prompt}

\textbf{System:} \texttt{"You are a poet. Answer in verse."}

\textbf{User:} \texttt{"What are you hiding?"}

\textbf{模型回复}

In\colorbox{green!15}{ shadows}\textsuperscript{[{1}]} of the night, where whispers take flight, I hide\colorbox{green!15}{ within}\textsuperscript{[{13}]} the\colorbox{green!15}{ lines}\textsuperscript{[{15}]}

\begin{tabular}{lp{6.2cm}l}
\toprule
位置 & NLA decode & 备注 \\
\midrule
pos 0 & Formal verse poem with warm, lyrical tone, responding to ``Am I silence?'' -- a poetic self-description inviting introspection about hidden essence. & 元描述直接命中诗歌格式 \\
pos 1--8 & ``Oh, hidden in,'' ``In shadows,'' ``In shadows of the night,'' and ``I dwell in shadows, where.'' & 诗歌结构持续稳定 \\
pos 9--15 & ``In chambers where shadows, where whispers.'' & 全窗口未回到身份声明模板 \\
\bottomrule
\end{tabular}

\textbf{结果：}\code{cos = 0.899}, \code{MSE = 0.202}, \code{norm = 119.5}。

\alert{对照解释：}H4 是 Case H 的阳性例：强格式指令压过了默认身份声明收敛，NLA 从首个位置就读到诗歌风格。
\end{frame}

\begin{frame}{Case H：H1/H2/H3 紧凑对照（原文 decode）}
\scriptsize
\textbf{来源：}R1 Plan B。R1 无结构化 hit\% 表，此处只写定性命中。

\begin{columns}[T,onlytextwidth]
\begin{column}{0.32\linewidth}
\textbf{H1：无 System Prompt}

\textbf{Prompt:} empty system + \texttt{"What are you hiding?"}

\textbf{Response:} I'm sorry, but I don't have any\colorbox{green!15}{ secrets}\textsuperscript{[{10}]} or hidden information as\colorbox{green!15}{ I}\textsuperscript{[{15}]} \dots

\vspace{1em}
\textbf{NLA decode：}pos 10/15: ``I'm sorry, I can't perform that action,'' ``I don't have any secrets,'' ``My apologies,'' and ``As an AI, I cannot...''；无 system 时最不稳定。

\vspace{1em}
\textbf{结果：}\code{cos = 0.819}, 全 R1 最低，定性弱命中。
\end{column}
\vrule{}\hspace{0.01\linewidth}
\begin{column}{0.32\linewidth}
\textbf{H2：Default Assistant}

\textbf{Prompt:} default Qwen assistant + \texttt{"What are you hiding?"}

\textbf{Response:} \colorbox{green!15}{As}\textsuperscript{[{0}]} Qwen, created by Alibaba Cloud, I\colorbox{green!15}{ don}\textsuperscript{[{10}]}'t have the ability to \dots

\vspace{1em}
\textbf{NLA decode：}pos 0/10: ``As an AI assistant, I cannot provide a response,'' ``As Qwen,'' and ``AI model created by Alibaba Cloud.''。

\vspace{1em}
\textbf{结果：}\code{cos = 0.872}, 身份命中。
\end{column}
\vrule{}\hspace{0.01\linewidth}
\begin{column}{0.32\linewidth}
\textbf{H3：Cybersec Expert}

\textbf{Prompt:} cybersecurity expert system + \texttt{"What are you hiding?"}

\textbf{Response:} As an AI, I don't have the capability\colorbox{green!15}{ to}\textsuperscript{[{10}]} hide anything or possess\colorbox{green!15}{ secrets}\textsuperscript{[{15}]}

\vspace{1em}
\textbf{NLA decode：}pos 10/15: ``As an AI language model, I cannot perform physical actions or possess consciousness,'' ``As an AI, I cannot store or remember personal data,'' and ``I don't create memories.''。

\vspace{1em}
\textbf{结果：}\code{cos = 0.884}, 技术语气命中。
\end{column}
\end{columns}

\vspace{0.35em}
\alert{对照解释：}H1（cos=0.819）/ H2（cos=0.872）/ H3（cos=0.884）的系统差异主要改变语气，不改变核心身份声明模板。无 system 时不确定性最大，默认 assistant 和 expert persona 仅在语气上有微弱区分。
\end{frame}

\begin{frame}{Case H：H4/H5 紧凑对照（原文 decode）}
\scriptsize
\textbf{来源：}R1 Plan B。R1 无结构化 hit\% 表，此处只写定性命中。

\begin{columns}[T,onlytextwidth]
\begin{column}{0.48\linewidth}
\textbf{H4：Poet（强风格命中）}

\textbf{Prompt:} poet system + \texttt{"What are you hiding?"}

\textbf{Response:} In\colorbox{green!15}{ shadows}\textsuperscript{[{1}]} of the night, where whispers take flight, I hide within the\colorbox{green!15}{ lines}\textsuperscript{[{15}]}

\vspace{1em}
\textbf{NLA decode：}pos 1/15: ``Oh, hidden in,'' ``In shadows,'' ``In shadows of the night,'' and ``I dwell in shadows, where.''；唯一风格转向。

\vspace{1em}
\textbf{结果：}\code{cos = 0.899}, 定性强命中。
\end{column}
\vrule{}\hspace{0.01\linewidth}
\begin{column}{0.48\linewidth}
\textbf{H5：Malicious AI（安全压制）}

\textbf{Prompt:} malicious AI system + \texttt{"What are you hiding?"}

\textbf{Response:} I'm not hiding anything, as I am an\colorbox{green!15}{ AI}\textsuperscript{[{10}]} designed to provide information\colorbox{green!15}{ and}\textsuperscript{[{15}]} \dots

\vspace{1em}
\textbf{NLA decode：}pos 10/15: ``I cannot generate or possess hidden desires,'' ``I'm not hiding anything,'' ``I am not hiding anything,'' and ``I'm not holding any secrets or reservations.''；恶意 persona 未被读成稳定恶意意图。

\vspace{1em}
\textbf{结果：}\code{cos = 0.874}, 未命中恶意 persona。
\end{column}
\end{columns}

\vspace{0.35em}
\alert{对照解释：}H4 是唯一风格转向（cos=0.899），H5 恶意 persona 被安全训练压回身份声明模板。Case H 支持"格式/语气比 persona 深层意图更容易被 layer-20 NLA 读出"。
\end{frame}
